\documentclass[12pt,a4paper]{article}
\usepackage[utf8]{inputenc}
\usepackage[T1]{fontenc}
\usepackage[polish]{polski}
\usepackage{amsmath}
\usepackage{amsfonts}
\usepackage{amssymb}
\usepackage{booktabs}
\usepackage{geometry}
\usepackage{enumitem}
\usepackage{graphicx}
\usepackage{float}
\usepackage{listings}
\usepackage{xcolor}
\usepackage{tabularx}
\usepackage{array}
\usepackage{etoolbox}

% lstlisting nigdy nie łamany między stronami
\BeforeBeginEnvironment{lstlisting}{\begin{minipage}{\linewidth}}
\AfterEndEnvironment{lstlisting}{\end{minipage}}

\geometry{margin=2cm}

\definecolor{codebg}{rgb}{0.97,0.97,0.97}
\definecolor{codekey}{rgb}{0.0,0.3,0.7}
\definecolor{codecomment}{rgb}{0.2,0.5,0.2}
\definecolor{codestring}{rgb}{0.6,0.1,0.1}

\lstset{
    language=C++,
    basicstyle=\ttfamily\small,
    keywordstyle=\color{codekey}\bfseries,
    commentstyle=\color{codecomment}\itshape,
    stringstyle=\color{codestring},
    backgroundcolor=\color{codebg},
    frame=single,
    framesep=4pt,
    breaklines=true,
    showstringspaces=false,
    numbers=left,
    numberstyle=\tiny\color{gray},
    xleftmargin=2em,
    tabsize=4,
    columns=fullflexible,
    literate=
        {ą}{{\k{a}}}1 {Ą}{{\k{A}}}1
        {ę}{{\k{e}}}1 {Ę}{{\k{E}}}1
        {ó}{{\'o}}1   {Ó}{{\'O}}1
        {ś}{{\'s}}1   {Ś}{{\'S}}1
        {ż}{{\.z}}1   {Ż}{{\.Z}}1
        {ź}{{\'z}}1   {Ź}{{\'Z}}1
        {ć}{{\'c}}1   {Ć}{{\'C}}1
        {ń}{{\'n}}1   {Ń}{{\'N}}1
        {ł}{{\l{}}}1  {Ł}{{\L{}}}1,
}

\title{Systemy czasu rzeczywistego, Projekt 4}
\author{Piotr Gugnowski, wydział EiTI, nr indeksu 320690}
\date{Czerwiec 2026}

\pagestyle{empty}

\begin{document}

\maketitle

% ═══════════════════════════════════════════════════════════════════════════════
\section*{Zadanie 4.1 — System CAN: suma kontrolna}
% ═══════════════════════════════════════════════════════════════════════════════

\subsection*{Treść zadania}

Jak wiemy, determinizm jest bardzo ważnym, ale nie jedynym i~często nie
najważniejszym wymaganiem stawianym sieciom czasu rzeczywistego. W~praktyce
ogromną rolę odgrywa bezpieczeństwo przesyłanych informacji. Pod tym pojęciem
rozumiemy w~węższym sensie zgodność informacji nadawanej i~odbieranej w~sieci.
W~szerszym sensie bezpieczeństwo informacji przesyłanych w~sieci polega również
na ich odpowiednim zabezpieczeniu w~celu utrudnienia lub uniemożliwienia dostępu
do nich stronom trzecim, lub/i~zastosowaniu odpowiednich technik kryptograficznych.
W~tym zadaniu interesować nas będzie wyłącznie bezpieczeństwo w~zdefiniowanym
wyżej węższym sensie. Jak wiadomo, na skutek zakłóceń komunikacyjnych np.\
wskutek oddziaływania silnych zaburzeń elektromagnetycznych, przesyłana
informacja może ulec zniekształceniu. Dla celów detekcji błędów transmisji
stosowane są różne techniki. Jedną z~nich jest stosowana dość powszechnie
w~profesjonalnych sieciach komunikacyjnych redundantna informacja zwana sumą
kontrolną lub cykliczną sumą redundancyjną. Suma kontrolna jest tworzona według
określonego algorytmu na podstawie wysyłanych informacji przez nadajnik.
Suma ta dołączana jest do strumienia wysyłanej informacji. Odbiornik tworzy własną
sumę kontrolną według tego samego algorytmu co nadajnik, ale na podstawie
strumienia danych odbieranych. Różnica w~odebranej i~zrekonstruowanej w~odbiorniku
sumie kontrolnej wskazuje na wystąpienie błędu lub błędów komunikacyjnych
i~pozwala na podjęcie odpowiednich akcji. Należy jednak pamiętać, że:

\begin{itemize}
    \item istnieje skończone prawdopodobieństwo zgodności sum kontrolnych odebranej
          i~utworzonej w~odbiorniku mimo wystąpienia błędu lub błędów transmisji;
    \item prawdopodobieństwo to jest w~profesjonalnych systemach sieciowych znikome
          i~minimalizowane dzięki odpowiedniemu doborowi algorytmów wyznaczania sumy
          kontrolnej i~starannemu doborowi wielomianów generacyjnych.
\end{itemize}

W~przypadku sieci CAN stosowany jest następujący algorytm.

\subsubsection*{Algorytm wyznaczania CRC}

Niech \texttt{CRC\_RG(14:0)} oznacza rejestr przechowujący wartość sumy kontrolnej CRC,
natomiast \texttt{NXTBIT} oznacza kolejny bit transmitowanego ciągu bitów począwszy od
SOF i~skończywszy na ostatnim bicie pola DATA. Suma CRC obliczana jest następująco:

\begin{lstlisting}[language={},basicstyle=\ttfamily\small,numbers=none,frame=single,backgroundcolor=\color{codebg}]
CRC_RG = 0;
REPEAT
    CRCNXT = NXTBIT EXOR CRC_RG(14)
    CRC_RG(14:1) = CRC_RG(13:0);
    CRC_RG(0) = 0;
    IF CRCNXT THEN
        CRC_RG(14:0) = CRC_RG(14:0) EXOR (4599hex);
    ENDIF
UNTIL (zostanie osiągnięte CRC SEQUENCE lub wystąpi błąd transmisji)
\end{lstlisting}

Napisać aplikację wyznaczającą sumę kontrolną dla sieci CAN. Aplikacja ta ma być
optymalizowana ze względu na minimalizację czasu realizacji algorytmu wyznaczania CRC.

\textbf{Założenia:}
\begin{enumerate}[label=\alph*)]
    \item Aplikacja ma umożliwiać wprowadzenie z~klawiatury ciągu do 96~bitów.
    \item Aplikacja powinna wyznaczyć sumę kontrolną z~wprowadzonego ciągu
          bitowego i~wyświetlić ją w~postaci heksadecymalnej.
    \item Aplikacja powinna mieć możliwość deklarowania liczby powtórzeń obliczenia
          sumy kontrolnej od~1 do~$10^9$.
    \item Aplikacja powinna mieć możliwość pomiaru i~wyświetlania łącznego czasu
          realizacji zadeklarowanej liczby powtórzeń obliczenia sumy kontrolnej oraz
          średniego czasu jednokrotnej realizacji obliczenia CRC.
    \item Aplikacja powinna być dostarczona w~postaci kodu maszynowego wykonywalnego.
\end{enumerate}

\subsection*{Implementacja}

Kod źródłowy znajduje się w~pliku \texttt{src/main.cpp}.

\subsubsection*{Reprezentacja wejścia}

Ciąg bitów wczytany z~klawiatury jest pakowany w~strukturę \texttt{BitInput}:
16 bajtów przechowuje bity MSB-first (pierwszy znak → bit~7 bajtu~0,
drugi → bit~6 bajtu~0 itd.), a~pole \texttt{nbits} przechowuje dokładną
liczbę bitów (0–96).

\subsubsection*{Tablica przeglądowa generowana w~czasie kompilacji}

Funkcja \texttt{can\_crc15()} oblicza CRC-15/CAN w~dwóch fazach:

\begin{enumerate}[label=\arabic*.]
    \item \textbf{Bajt po bajcie} (dla $\lfloor n/8 \rfloor$ pełnych bajtów)
          przy użyciu 256-elementowej tablicy \texttt{TABLE} (512~B, mieści się
          w~L1~cache). Klucz do tablicy: \texttt{key = ((crc >> 7) \^{} byte) \& 0xFF};
          aktualizacja: \texttt{crc = ((crc << 8) \& 0x7FFF) \^{} TABLE[key]}.
          Daje 8-krotną redukcję liczby iteracji względem przetwarzania
          bit po bicie.
    \item \textbf{Pozostałe bity} (0–7, ostatni niepełny bajt): iteracja
          bit po bicie z~bezgałęziowym krokiem \\
          \texttt{crc \^{}= crcnxt * 0x4599u} (eliminacja branch misprediction).
\end{enumerate}

Tablica jest obliczana jako \texttt{constexpr}.

\subsubsection*{Zabezpieczenia integralności benchmarku}

\begin{itemize}
    \item \textbf{\texttt{volatile uint16\_t sink}} — kompilator nie może
          wynieść wywołania funkcji poza pętlę (efekt uboczny obserwowalny
          z~zewnątrz).
    \item \textbf{\texttt{\_\_attribute\_\_((noinline))}} na funkcji
          \texttt{can\_crc15} — blokuje inlining, przez co kompilator nie
          widzi wnętrza funkcji i~nie może udowodnić stałości wyniku.
\end{itemize}

\subsubsection*{Pełny kod funkcji CRC}

\begin{lstlisting}
uint16_t can_crc15(const BitInput& in) noexcept {
    uint16_t crc = 0;
    const int full = in.nbits / 8;
    const int rem  = in.nbits % 8;

    for (int i = 0; i < full; ++i) {
        const uint8_t key =
            static_cast<uint8_t>(((crc >> 7) ^ in.bytes[i]) & 0xFFu);
        crc = static_cast<uint16_t>(
            ((crc << 8) & 0x7FFFu) ^ TABLE.v[key]);
    }

    if (rem) {
        const uint8_t last = in.bytes[full];
        for (int i = 7; i >= 8 - rem; --i) {
            const uint8_t crcnxt = static_cast<uint8_t>(
                ((crc >> 14) ^ ((last >> i) & 1u)) & 1u);
            crc = static_cast<uint16_t>((crc << 1) & 0x7FFFu);
            crc ^= static_cast<uint16_t>(crcnxt * 0x4599u);
        }
    }
    return crc;
}
\end{lstlisting}

\subsection*{Wyniki i~pomiary czasu}

Pomiary wykonano na Apple Silicon (ARM64), kompilacja:
\texttt{clang++ -O3 -march=native}, $n = 10^9$ powtórzeń.

\begin{table}[H]
\centering
\begin{tabular}{llrrr}
\toprule
\textbf{Wejście} & \textbf{Bity} & \textbf{CRC-15/CAN} & \textbf{Czas łączny} & \textbf{Czas/iter.} \\
\midrule
96 zer  (\texttt{00...0}) & 96 & \texttt{0x0000} & 319,5~ms & 0,320~ns \\
96 jedynek (\texttt{11...1}) & 96 & \texttt{0x072A} & 313,7~ms & 0,314~ns \\
\bottomrule
\end{tabular}
\caption{Wyniki benchmarku CRC-15/CAN dla $n = 10^9$ iteracji.}
\end{table}

Wynik dla 96~zer (\texttt{0x0000}) jest zgodny z~oczekiwaniami:
wszystkie bity wejściowe zerowe powodują, że \texttt{CRCNXT} jest zawsze
zerem, rejestr CRC pozostaje zerowy przez cały czas obliczeń.

% ═══════════════════════════════════════════════════════════════════════════════
\section*{Zadanie 4.2 — System CAN: identyfikatory wiadomości}
% ═══════════════════════════════════════════════════════════════════════════════

\subsection*{Treść zadania}

Jak wiadomo, węzły CAN są generatorami wiadomości. Każda wiadomość ma swój
własny identyfikator. Zależnie od specyfikacji sieci CAN występują w~specyfikacji
standardowej identyfikatory 11-bitowe, a~w~specyfikacji rozszerzonej identyfikatory
29-bitowe. W~związku z~tym teoretycznie możliwe jest zaadresowanie maksymalnie
odpowiednio $2^{11}$ lub $2^{29}$ wiadomości.

Czy tak jest w~rzeczywistości? Proszę podać istniejące ograniczenia w~specyfikacji
CAN (o~ile występują).

\subsection*{Odpowiedź}

Nie, nie wszystkie kombinacje identyfikatorów są dostępne. Specyfikacja CAN zabrania
używania identyfikatorów, których \textbf{pierwsze 7~bitów} to same jedynki
(\texttt{1111111\dots}). Taka kombinacja wygląda na magistrali identycznie jak sygnał
błędu pasywnego (\textit{passive error flag} - 6~bitów recessive) lub jak przerwa
między ramkami (IFS), co mogłoby wprowadzić odbiorniki w~błąd.

\begin{table}[H]
\centering
\begin{tabular}{lrrr}
\toprule
\textbf{Format} & \textbf{Bity ID} & \textbf{Zakazane ID} & \textbf{Dostępne ID} \\
\midrule
Standardowy & 11 & $2^{4} = 16$ & $2048 - 16 = \mathbf{2032}$ \\
Rozszerzony & 29 & $2^{22} = 4\,194\,304$ &
    $536\,870\,912 - 4\,194\,304 = \mathbf{532\,676\,608}$ \\
\bottomrule
\end{tabular}
\caption{Rzeczywista liczba dostępnych identyfikatorów w~sieci CAN.}
\end{table}

W~praktyce ograniczenie to jest nieistotne, typowa sieć przemysłowa korzysta
z~kilkudziesięciu do kilkuset typów wiadomości, co stanowi ułamek procenta
dostępnej przestrzeni adresowej.

% ═══════════════════════════════════════════════════════════════════════════════
\section*{Zadanie 4.3 — System CAN: rzeczywisty czas trwania transmisji ramki}
% ═══════════════════════════════════════════════════════════════════════════════

\subsection*{Treść zadania}

Dokładna analiza czasu transmisji ramki CAN wskazuje, że przy transmisji tej samej
liczby bajtów informacji w~polu danych czas trwania transmisji jest różny. Czas ten
zależy od kontekstu. Przyczyna leży w~sposobie kodowania informacji, a~w~szczególności
w~zastosowanej technice synchronizacji transferów sieciowych wspieranych przez
dodatkowe bity synchronizujące (technika \textit{bit stuffing}).

\textbf{Założenia:}
\begin{enumerate}[label=\alph*)]
    \item Prędkość transmisji w~sieci CAN: 1\,Mb/s.
    \item Przesyłamy w~polu danych 8~bajtów informacji.
    \item Identyfikator wiadomości podlega swobodnemu wyborowi z~ew.\ ograniczeniami
          (por.\ zadanie~4.2).
    \item Stosujemy ramkę o~identyfikatorze 29-bitowym.
    \item Współczynnik efektywności transmisji zdefiniowany jako stosunek liczby bitów
          wysłanej informacji (64~bity) do całkowitej liczby bitów ramki (łącznie
          z~bitami przerwy między ramkami i~bitami dodatkowymi dodanymi przez operację
          \textit{bit stuffing}).
\end{enumerate}

\begin{enumerate}[label=\alph*)]
    \item Wyznaczyć najkrótszy czas transmisji całej ramki CAN. Podać przykład
          zawartości pola danych, dla której ma to miejsce. Podać wartość współczynnika
          efektywności transmisji.
    \item Wyznaczyć najdłuższy czas transmisji całej ramki CAN. Podać przykład
          zawartości pola danych, dla której ma to miejsce. Podać wartość współczynnika
          efektywności transmisji.
\end{enumerate}

\subsection*{Wyprowadzenie liczby bitów ramki}

Rozszerzona ramka CAN~2.0B z~8~bajtami danych składa się z~następujących pól:

\begin{table}[H]
\centering
\small
\begin{tabular}{lr}
\toprule
\textbf{Pole} & \textbf{Bitów} \\
\midrule
SOF                    &  1 \\
Arbitration (ID+flagi) & 32 \\
Control                &  6 \\
Data                   & 64 \\
CRC sequence           & 15 \\
CRC delimiter          &  1 \\
ACK slot               &  1 \\
ACK delimiter          &  1 \\
EOF                    &  7 \\
IFS                    &  3 \\
\midrule
\textbf{Łącznie}       & \textbf{131} \\
\bottomrule
\end{tabular}
\end{table}

\begin{equation*}
L_{\text{stały}} = 1 + 32 + 6 + 64 + 15 + 1 + 1 + 1 + 7 + 3 = \mathbf{131\ \text{bitów}}
\end{equation*}

Technika \textbf{bit stuffing} polega na automatycznym wstawianiu przez nadajnik
jednego bitu o~przeciwnej wartości po każdych 5~kolejnych bitach tej samej wartości.
Odbiornik usuwa te dodatkowe bity. Cel: utrzymanie synchronizacji zegara odbiornika,
który nie ma osobnego sygnału taktowania i~synchronizuje się na zmianach stanu
na magistrali.

Bit stuffing obowiązuje dla regionu: SOF + Arbitration + Control + Data + CRC sequence:
\begin{equation*}
N_{\text{fill}} = 1 + 32 + 6 + 64 + 15 = \mathbf{118\ \text{bitów}}
\end{equation*}

Przy prędkości 1\,Mb/s czas jednego bitu wynosi $T_b = 1\,\mu$s.

% ─────────────────────────────────────────────────────────────────────────────
\subsection*{a) Najkrótszy czas transmisji}

Najkrótszy czas uzyskuje się przy \textbf{zerowej liczbie bitów stuffujących},
co zachodzi gdy w~regionie 118~bitów nie pojawi się żadna sekwencja 5~kolejnych
bitów tej samej wartości.

\textbf{Przykład danych spełniających ten warunek:}
\begin{center}
\texttt{0x55 0x55 0x55 0x55 0x55 0x55 0x55 0x55}
\end{center}
Bajt \texttt{0x55} = \texttt{01010101}$_2$ — bity naprzemienne, nie tworzą
żadnej sekwencji 5~identycznych.
Identyfikator musi mieć analogiczną właściwość: żaden fragment jego 29~bitów
nie może tworzyć serii 5~identycznych, a~ponadto złożenie SOF\,+\,ID\,+\,pola
Control (łącznie z~wynikową sumą CRC) nie może stworzyć takiej serii na granicy
sąsiednich pól. W~praktyce identyfikator z~naprzemiennymi bitami
(postaci \texttt{01010\ldots} lub \texttt{10101\ldots}) spełnia ten warunek
przy naprzemiennych danych \texttt{0x55}.

\begin{align*}
L_{\min} &= L_{\text{stały}} + N_{\text{stuff}} = 131 + 0 = \mathbf{131\ \text{bitów}}
\end{align*}
\[
\boxed{T_{\min} = 131 \cdot 1\,\mu\text{s} = \mathbf{131\,\mu\text{s}}}
\]

\textbf{Współczynnik efektywności transmisji} (stosunek bitów użytecznych do
wszystkich bitów ramki):
\[
\eta_{\max} = \frac{64}{131} \approx \mathbf{48{,}9\%}
\]

% ─────────────────────────────────────────────────────────────────────────────
\subsection*{b) Najdłuższy czas transmisji}

Najdłuższy czas uzyskuje się przy \textbf{maksymalnej liczbie bitów stuffujących}.
Bit stuffujący jest wstawiany po każdych 5~bitach tej samej wartości, więc
najgorszy wzorzec to grupy pięciu identycznych bitów:
\texttt{11111\,00000\,11111\,00000}\ldots

Dla regionu 118~bitów każda grupa 5~bitów wymusi wstawienie 1~bitu stuffującego:
\[
N_{\text{stuff,max}} = \left\lfloor \frac{118}{5} \right\rfloor = \mathbf{23}
\]

\textbf{Przykład danych generujących dużo stuffowania:}
\begin{center}
\texttt{0x00 0x00 0x00 0x00 0x00 0x00 0x00 0x00}
\end{center}
64~kolejne zera w~polu danych samo w~sobie daje
$\lfloor 64/5 \rfloor = 12$ bitów stuffujących.
Aby osiągnąć 23~bity stuffujące łącznie, identyfikator powinien zawierać
długie serie identycznych bitów (np.\ postaci \texttt{11111\,00000\ldots}),
a~wynikowa suma kontrolna CRC — obliczana automatycznie z~ID i~danych,
bez możliwości swobodnego wyboru — musi wspólnie z~ID wypełnić pozostałe
$23 - 12 = 11$ bitów stuffujących w~regionie poza polem danych.
Znalezienie wartości ID spełniającej ten warunek wymaga iteracyjnego
sprawdzania kandydatów z~uwzględnieniem wynikowego CRC.

\begin{align*}
L_{\max} &= L_{\text{stały}} + N_{\text{stuff,max}} = 131 + 23 = \mathbf{154\ \text{bitów}}
\end{align*}
\[
\boxed{T_{\max} = 154 \cdot 1\,\mu\text{s} = \mathbf{154\,\mu\text{s}}}
\]

\textbf{Współczynnik efektywności transmisji:}
\[
\eta_{\min} = \frac{64}{154} \approx \mathbf{41{,}6\%}
\]

% ─────────────────────────────────────────────────────────────────────────────
\subsection*{Podsumowanie}

\begin{center}
\begin{tabular}{llllrrrr}
\toprule
\textbf{Scenariusz} & \textbf{Dane (8 B)} & \textbf{Bity stałe} & \textbf{Bity stuff.} & \textbf{Łącznie} & \textbf{Wsp. ef. trans. $\bm{\eta}$} \\
\midrule
Najkrótsza ramka & \texttt{0x55}$\times$8 & 131 & 0  & 131 & 48,9\% \\
Najdłuższa ramka & \texttt{0x00}$\times$8 & 131 & 23 & 154 & 41,6\% \\
\bottomrule
\end{tabular}
\end{center}


% Tło teoretyczne i szczegóły implementacyjne — patrz teoria.pdf

\end{document}
